上一节中已经在热力学极限态及其两个两格点转移算符 $E_L$ 和 $E_R$ 中发现了同样的结构。这使得旧有结果可以直接移植过来。假设我们想要计算 $\bra{\psi} \hat{O}_1 \hat{O}_i \ket{\psi}$；那么我们把该态化成式~(\ref{eq:infinitemixedcanonical}) 的混合规范形式，即 $A$ 矩阵一直排到格点 $i$ 或 $i+1$（取决于奇偶结构），然后是 $\Lambda^{[\ell-1]} $，再往后是直到无穷的 $B$ 矩阵。从左边收缩掉所有直到 0 的 $A$ 矩阵、从右边收缩掉所有 $B$ 矩阵之后，我们就得到如图~\ref{fig:finitecontractionremainder} 所示的剩余有限网络。
这个期望值随后可以像有限网络那样求值；如果我们使用 $C^{[i]}$ 与转移算符的记法，则从 $C^{[0]}=I$ 出发。区别在于，在最后 $\Lambda^{[\ell-1]}$ 会出现在收缩之中，而最终约化为标量则由收尾的 $\delta_{aa'}$ 线完成（它可以被理解为求迹）：假设最后一个格点是 $i$，那么最终的期望值在最后一步由下式给出
\begin{equation}
\bra{\psi} O^1 \otimes \ldots O^i \ket{\psi} = \tr  \Lambda^{[\ell-1]\dagger} C^{[i]} \Lambda^{[\ell-1]} = 
\tr   \Lambda^{[\ell-1]} \Lambda^{[\ell-1]\dagger} C^{[i]} = \tr \rho_A C^{[i]} ,
\end{equation}
其中我们用到了 $\Lambda$ 矩阵与规范形式下约化密度算符之间的关系。

\begin{figure}
\centering\includegraphics[scale=0.7]{PyMPS_FiniteContractionRemainder.eps}
\caption{利用左归一化和右归一化，无限收缩的求值可以约化为一个有限收缩，从而可以像处理任何有限 MPS 那样处理它。}
\label{fig:finitecontractionremainder}
\end{figure}
    
这一计算可以很容易地约化到两个态的重叠的情形，后者不过是单位算符在这两个态之间的矩阵元。假设两个态都处于平移不变形式，例如分别使用左归一化矩阵 $A^{[1]}$ 和 $A^{[2]}$（以及 $\tilde{A}^{[1]}$、$\tilde{A}^{[2]}$）。现在我们用 $E_L$ 把一个无限重叠计算向右推进两个格点（比如 1 和 2）：如果当前的重叠矩阵是 $C$，则它按如下方式推进：
\begin{equation}
E_L (C) = \sum_{\sigma_1\sigma_2} \tilde{A}^{[2]\sigma_2\dagger} \tilde{A}^{[1]\sigma_1\dagger} C A^{[1[\sigma_1} A^{[2]\sigma_2} .
\end{equation}  
如果把 $C$ 按 $E_L$ 的本征矩阵分解，那么在热力学极限下只有最大本征值的贡献会存活下来。对于一个正交归一的态，它与自身的重叠给出主导本征对 $C=I$ 和 $\lambda=1$。两个态的重叠中较小的 $\lambda$ 可以被解读为每格点重叠；当然，在热力学极限下这两个态彼此是正交的（重叠 $\lim_{L\rightarrow\infty} \lambda^L = 0$）。这种每格点的热力学重叠（或保真度）可以非常好地用来探测量子相变：只需把参数稍有不同的两个哈密顿量的基态作重叠\cite{McCulloch08}。
 
\section{结论：其他值得关注的课题}
在遍历了这份长长的课题清单——聚焦于基本的算法构件、基态搜索、热力学极限算法以及零温和有限温下丰富的实时与虚时方法——之后，可以看到 DMRG 与基于 MPS 的算法的可能性还远未穷尽。有几个我未曾论及、但想简要提及的课题（同样是一份不完全的清单）是：转移矩阵 DMRG 方法（参考文献见引言部分）、带周期边界条件的 DMRG 与 MPS，以及作为最新进展的连续空间 MPS\cite{Verstraete10}——它表现为相干态的一种漂亮推广，应当能在场论中找到有趣的应用。对于 {\em 周期边界条件}，已有的结果已经相当多，所以这里只做一个简短的概述。在 DMRG 框架内处理 PBC 的做法，是在一条开边界链的格点 1 与 $L$ 之间引入一个长程相互作用（参见例如 \cite{Schmitteckert96,Rapsch99,Meden03a,Meden03b}；然而人们一致发现其精度标度比开边界条件差得多。其深层原因（粗略地说）在于：在环上 A 与 B 之间的界面加倍，纠缠也随之加倍；考虑到纠缠与 MPS 维数之间的指数关系，这意味着资源必须从 $D$ 增至最多 $D^2$，也就是说，为达到相近的精度，算法需要平方的时间（有时被称为 $D^6$ 标度，这是相对于开边界条件的 $D$ 而言的）。物理上恰当的周期边界条件 MPS 拟设由式 (\ref{eq:MPSforPBC}) 给出；所需的 $D$ 与 OBC 大致相同，但在其上重新运行变分基态搜索算法的标度是 $D^5$（因为在格点 1 和 $L$ 上矢量代替矩阵的简化不再出现）\cite{VerstraetePorras04}。同时，从广义本征值问题到标准本征值问题的简化也不再发生，这可能导致病态的条件数。PBC 的 MPS 表示有一个很好的特点：可以生成动量的本征态。对于 $k = n (2\pi/L)$ 以及一个（非平移不变的）MPS $\ket{\psi} = \sum_{{\fat{\sigma}}} \tr (A^{[1]\sigma_1} \ldots A^{[L]\sigma_L}) \ket{\fat{\sigma}}$，下面的态就是动量 $k$ 的平移不变本征态：\cite{Porras06}
\begin{equation}
\ket{\psi_k} = \sum_{n=0}^{L-1} \eul^{\imag k n}  \tr (A^{[1]\sigma_{1+n}} \ldots A^{[L]\sigma_{L+n}}) \ket{\fat{\sigma}} .
\end{equation}
最近已有一些改进 $D^5$ 标度的有趣方案被提出 \cite{Pippan10,Pirvu10}，这是一个仍在持续受到关注的领域。文献 \cite{VerstraeteMurg08} 相当详尽地讨论了这一课题。

我想可以这样作结：虽然 MPS 的基本框架如今已经非常成熟，DMRG 也已成长为强关联量子体系可用的最强大数值方法之一，但即使在一维体系这一发展完善的领域中，本文介绍的许多算法仍有进一步改进的余地，从而把新的应用带入我们的视野。事实上相当令人惊讶的是，对于这里介绍的不少方法（其他方法也一样），人们对它们在真实问题中的细致行为知之甚少，而分析这些行为或许会带来有趣的进一步想法。此外，已完成的应用与可完成的应用之比，对未来的激动人心的研究而言似乎十分有利。

\section{致谢}
我要感谢 Aspen Center for Physics、意大利的里雅斯特的 International Center for Theoretical Physics 以及柏林的 Wissenschaftskolleg（Institute for Advanced Study）的热情接待，本工作的主要部分就是在这些地方完成的。兼具激发思考的讨论与潜心钻研的空间，正是物理学家梦寐以求的环境。当然还要感谢许多与我讨论过这些课题、并使我受益良多的人：Mari-Carmen Banuls, Thomas Barthel, Ignacio Cirac, Andrew Daley, Jan von Delft, Jens Eisert, Adrian Feiguin, Juanjo Garcia-Ripoll, Fabian Heidrich-Meisner, Corinna Kollath, Ian McCulloch, Valentin Murg, Volker Nebendahl, Tomotoshi Nishino, Reinhard Noack, Tobias Osborne, Lode Pollet, Matthias Troyer, Frank Verstraete, Guifre Vidal, Andreas Weichselbaum, Steve White 以及 Tao Xiang，此处仅举其要者。特别感谢 Matt Hastings 向我指出他的一篇重要论文，并感谢 Stefan Depenbrock 对本综述一个接近定稿的版本所作的极为仔细的阅读（所有残留的错误都是在那之后加入的）。

\begin{thebibliography}{999}

\bibitem{Bloch08}
I.~Bloch, J.~Dalibard and W.~Zwerger, Rev.\ Mod.\ Phys.\ {\bf 80}, 885 (2008).

\bibitem{Bethe31}
H.~Bethe, Z.\ Phys.\ A {\bf 71}, 205 (1931).

\bibitem{Lieb68}
E.~H.~Lieb and F.~Y.~Wu, Phys.\ Rev.\ Lett.\ {\bf 20}, 1445 (1968).

\bibitem{Haldane83}
F.~D.~M.~Haldane, Phys.\ Rev.\ Lett.\ {\bf 50}, 1153 (1983). 

\bibitem{White92}
S.~R. White, Phys.\ Rev.\ Lett.\ {\bf 69}, 2863 (1992). 

\bibitem{White93}
S.~R. White, Phys.\ Rev.\ B {\bf 48}, 10345 (1993). 

\bibitem{Schollwoeck05} U.~Schollw\"{o}ck, Rev.\ Mod.\ Phys. {\bf 77}, 259 (2005).

\bibitem{Hallberg06} K.~Hallberg, Adv.\ Phys.\ {\bf 55}, 477 (2006).

\bibitem{Hallberg95}
K.~Hallberg, Phys.\ Rev.\ B {\bf 52}, 9827 (1995).

\bibitem{Ramasesha97}
S.~Ramasesha, S.~K.~Pati, H.~R.~Krishnamurthy, Z.~Shuai and J.~L.~Br{\'e}das, Synth.\ Met.\ {\bf 85}, 1019 (1997).

\bibitem{Kuhner99}
T.~D. K\"{u}hner and S.~R. White, Phys.\ Rev.\ B {\bf 60}, 335 (1999).

\bibitem{Jeckelmann02}
E.~Jeckelmann, Phys.\ Rev.\ B {\bf 66}, 045114 (2002).
 
\bibitem{Nishino95}
T.~Nishino, J. Phys.\ Soc.\ Jpn.\ {\bf 64}, 3598 (1995).

\bibitem{Wang97}
X.~Q. Wang and T.~Xiang, Phys.\ Rev.\ B {\bf 56}, 5061 (1997).

\bibitem{Shibata97}
N.~Shibata, J. Phys.\ Soc.\ Jpn. {\bf 66}, 2221 (1997).

\bibitem{Sirker05}
J.~Sirker and A.~Kl\"{u}mper, Phys.\ Rev.\ B {\bf 71}, 241101 (2005).

\bibitem{Hieida98}
Y.~Hieida, J. Phys.\ Soc.\ Jpn. {\bf 67}, 369 (1998).

\bibitem{Carlon99}
E.~Carlon, M.~Henkel and U.~Schollw\"{o}ck, Eur.\ J. Phys.\ B {\bf 12}, 99 (1999).

\bibitem{Carlon01}
E.~Carlon, M.~Henkel and U.~Schollw\"{o}ck, Phys.\ Rev.\ E {\bf 63}, 036101 (2001).

\bibitem{Henkel01}
M.~Henkel and U.~Schollw\"{o}ck, J.\ Phys.\ A {\bf 34}, 3333 (2001).

\bibitem{White96}
S.~R.~White, Phys.\ Rev.\ Lett. {\bf 77}, 3633 (1996).

\bibitem{White98a} 
S.~R.~White and D.~J.~Scalapino, Phys.\ Rev.\ Lett.\ {\bf 80}, 1272 (1998).

\bibitem{White98b} 
S.~R.~White and D.~J.~Scalapino, Phys.\ Rev.\ Lett.\ {\bf 81}, 3227 (1998).

\bibitem{White00} 
S.~R.~White and D.~J.~Scalapino, Phys.\ Rev.\ B {\bf 61}, 6320 (2000).

\bibitem{White03} 
S.~R.~White and D.~J.~Scalapino, Phys.\ Rev.\ Lett.\ {\bf 91}, 136403 (2003).

\bibitem{Chernyshev03}
A.~L.~Chernyshev, S.~R.~White and A.~H.~Castro-Neto, Phys.\ Rev.\ B {\bf 65}, 214527 (2003).

